\section{Extensions: applications}
\begin{frame}[plain]{Extensions: empirical applications}
\Large Applications after the theory
\vspace{1em}\normalsize
\begin{enumerate}\setlength{\itemsep}{.65em}
\item Senate elections: local randomization
\item Ser Pilo Paga: imperfect compliance
\item Academic probation: a discrete score
\item Ser Pilo Paga and media markets: multiple dimensions
\end{enumerate}
\src{Stata modules 2--6 and worksheets B--F. The sequence follows Extensions.}
\end{frame}

\begin{frame}{Senate elections and party advantage}
\begin{items}
\item \textbf{Units:} 1,390 U.S. Senate election observations in the teaching dataset.
\item \textbf{Score:} Democratic margin of victory at election $t$.
\item \textbf{Treatment:} a Democratic win at $t$.
\item \textbf{Outcome:} Democratic vote share at election $t+2$, the next election for that same seat.
\item Here $t+2$ indexes elections. It does not mean two calendar years later.
\end{items}
\src{\textcite{cft2015}; Extensions, Section 2.1.}
\end{frame}

\begin{frame}{The Senate outcome-score relationship}
\fig[.73]{senate_plot.pdf}
\src{Module 2. Separate linear descriptive fits and automatically chosen bins. Outcome and score use percentage units.}
\end{frame}

\begin{frame}[fragile]{Local randomization in Stata}
\begin{lstlisting}
local covs presdemvoteshlag1 demvoteshlag1 demvoteshlag2 ///
    demwinprv1 demwinprv2 dmidterm dpresdem dopen
rdwinselect X `covs', wobs(2) nwindows(50) approx seed(50)
local wl=r(w_left)
local wr=r(w_right)
rdrandinf Y X, wl(`wl') wr(`wr') reps(10000) seed(50)
\end{lstlisting}
\begin{items}
\item The selected window is approximately $[-0.7652,0.7652]$ pp.
\item Balance tests choose the window. The outcome does not enter that selection.
\end{items}
\EX{Section 2.2; module 2}
\end{frame}

\begin{frame}{Cutoff and window estimates in the Senate data}
\begin{center}
\small\renewcommand{\arraystretch}{1.3}
\begin{tabular}{@{}lrrr@{}}\toprule
 & Estimate & Bias-corrected & Robust 95\% CI\\\midrule
Continuity at cutoff & 7.41 & 7.51 & $[4.09, 10.92]$\\
\bottomrule\end{tabular}
\end{center}

\begin{items}
\item The selected-window difference in means is $10.20$ percentage points.
\item None of 10,000 simulated null assignments gives a statistic as extreme as observed.
\item The Fisherian test concerns the sharp null of zero effect for every unit in the window.
\item The cutoff estimate and window estimate have different target populations and assumptions.
\end{items}
\src{Module 2. Continuity table: percentage points, robust bias corrected interval.}
\end{frame}

\begin{frame}{Ser Pilo Paga and higher education access}
\begin{items}
\item Colombia's programme combines academic merit and economic need.
\item The one-score sample contains 23,132 students who meet the academic requirement.
\item $X1$ measures poverty eligibility with its sign oriented so $X1\geq0$ means eligible.
\item $D$ records subsidy receipt. $Y$ records immediate enrolment in any higher education institution.
\item Eligible students need not take up the subsidy, so receipt follows a fuzzy design.
\end{items}
\src{\textcite{londono2020}; Extensions, Section 3.1.}
\end{frame}

\begin{frame}{The SPP first stage}
\fig[.73]{spp_D.pdf}
\src{Module 3. The vertical scale is the proportion receiving the subsidy, not a percentage.}
\end{frame}

\begin{frame}{The SPP reduced form}
\fig[.73]{spp_Y.pdf}
\src{Module 3. The outcome is immediate higher education enrolment. The discontinuity estimates an eligibility effect.}
\end{frame}

\begin{frame}{Eligibility, take-up, and the receipt effect}
\begin{center}
\small\renewcommand{\arraystretch}{1.3}
\begin{tabular}{@{}lrrr@{}}\toprule
 & Estimate & Bias-corrected & Robust 95\% CI\\\midrule
First stage (pp) & 62.48 & 62.35 & $[59.52, 65.19]$\\
Eligibility effect (pp) & 26.90 & 27.41 & $[22.06, 32.75]$\\
Receipt LATE (pp) & 43.45 & 44.51 & $[36.60, 52.42]$\\
\bottomrule\end{tabular}
\end{center}

\begin{items}
\item All effects and intervals in the table use percentage points.
\item The fuzzy estimate is a receipt LATE for cutoff compliers under the stated assumptions.
\item The first-stage row uses its own optimal bandwidth. Dividing the first two rows will not reproduce the joint fuzzy estimate.
\end{items}
\src{Module 3. A common-window calculation in the do-file reproduces the conventional Wald ratio.}
\end{frame}

\begin{frame}{A window selection failure in the SPP data}
\begin{items}
\item The module's approximate covariate balance procedure does not accept the smallest window at its stated threshold.
\item It therefore selects no local randomization window.
\item The code also shows the book's fixed $[-0.13,0.13]$ window as an explicitly labelled illustration.
\item That calculation does not establish a credible local randomization analysis under the failed selection protocol.
\end{items}
\src{Module 3; Extensions, Sections 3.2--3.3. The teaching conclusion retains the failed diagnostic.}
\end{frame}

\begin{frame}{Academic probation and a discrete GPA score}
\begin{items}
\item The teaching dataset contains 40,582 students and 429 distinct score values.
\item The original rule assigns probation below a campus-specific GPA standard.
\item The supplied transformed score is positive for probation and negative for control.
\item Exact original ties are controls. The authors encode them as $X=-0.000005$.
\item Next-term GPA is observed only for students who continue at university, creating a potential outcome-selection problem.
\end{items}
\src{\textcite{lindo2010}; Extensions, Sections 4.1--4.2.}
\end{frame}

\begin{frame}{The mass points in the probation score}
\fig[.73]{probation_histogram.pdf}
\src{Module 4. The large overall sample does not remove the discrete support. Closest groups: 208 controls at $-0.000005$, 67 treated at $0.01$.}
\end{frame}

\begin{frame}{The probation outcome plot}
\fig[.73]{probation_plot.pdf}
\src{Module 4. The score orientation makes lower prior performance correspond to a higher transformed score.}
\end{frame}

\begin{frame}{The probation estimates and their limits}
\begin{center}
\small\renewcommand{\arraystretch}{1.3}
\begin{tabular}{@{}lrrr@{}}\toprule
 & Estimate & Bias-corrected & Robust 95\% CI\\\midrule
Continuity method & 0.22 & 0.22 & $[0.13, 0.30]$\\
Cluster by score & 0.22 & 0.21 & $[0.14, 0.28]$\\
\bottomrule\end{tabular}
\end{center}

\begin{items}
\item Outcome units: GPA points. A point estimate near $0.22$ is not 22 percentage points.
\item In the closest two-score window, the difference is $0.234$ GPA points and the simulated Fisherian $p$-value is about $0.057$.
\item Assignment counts and covariate diagnostics raise credibility questions. Clustering by score does not settle them.
\end{items}
\src{Module 4. GPA among observed continuers does not automatically identify a causal effect for a common survivor population.}
\end{frame}

\begin{frame}{SPP with three area-specific cutoffs}
\begin{items}
\item The merit-eligible sample faces different poverty thresholds by area.
\item The supplied score reverses the original wealth-index sign.
\item Metropolitan areas use $c=-57.21$, other urban areas $c=-56.32$, and rural areas $c=-40.75$.
\item The outcome is enrolment, and the discontinuity is an eligibility effect.
\item This dataset does not include receipt, so these estimates do not by themselves identify subsidy receipt LATEs.
\end{items}
\EX{Section 5.1.3; module 5}
\end{frame}

\begin{frame}{Cutoff-specific and pooled SPP effects}
\begin{center}
\small\renewcommand{\arraystretch}{1.3}
\begin{tabular}{@{}lrrr@{}}\toprule
 & Estimate & Bias-corrected & Robust 95\% CI\\\midrule
Metropolitan (-57.21) & 34.64 & 36.02 & $[26.86, 45.18]$\\
Other urban (-56.32) & 20.29 & 19.69 & $[11.16, 28.23]$\\
Rural (-40.75) & 20.89 & 18.27 & $[-4.24, 40.79]$\\
Normalized pool & 26.90 & 27.41 & $[22.06, 32.75]$\\
\bottomrule\end{tabular}
\end{center}

\begin{items}
\item Units: percentage points of immediate enrolment.
\item Rural estimates are substantially less precise.
\item The normalized pooled estimate can conceal differences across groups.
\end{items}
\src{Module 5. The do-file also reports rdmc's estimated aggregation weights.}
\end{frame}

\begin{frame}{SPP eligibility with both scores}
\fig[.73]{spp_two_scores.pdf}
\src{Module 6, reproducible 2\% plotting sample. Eligibility requires both centered scores to be nonnegative. Estimation uses the full supplied dataset.}
\end{frame}

\begin{frame}{SPP effects at different boundary points}
\begin{center}
\small\renewcommand{\arraystretch}{1.3}
\begin{tabular}{@{}lrrr@{}}\toprule
 & Estimate & Bias-corrected & Robust 95\% CI\\\midrule
(0, 0) & 32.31 & 33.59 & $[29.30, 37.88]$\\
(30, 0) & 31.52 & 32.08 & $[28.58, 35.59]$\\
(0, 50) & 22.89 & 24.74 & $[14.39, 35.08]$\\
\bottomrule\end{tabular}
\end{center}

\begin{items}
\item Coordinates are (poverty eligibility score, exam score), both centered.
\item The data contain 574,269 student observations before outcome and score exclusions.
\item Each row uses signed distance to its boundary point, in the authors' score units.
\item These are eligibility effects in percentage points and need not coincide across locations.
\end{items}
\EX{Section 5.2.2; module 6}
\end{frame}

\begin{frame}{Political advertising at a media-market boundary}
\fig[.73]{geographic_boundary.pdf}
\src{Module 6; \textcite{keele2015}. Registered voters in New Jersey's West Windsor-Plainsboro school district. Outcome: turnout in the 2008 election.}
\end{frame}

\begin{frame}{Local support near the geographic boundary}
\fig[.73]{geographic_distance.pdf}
\src{Module 6. Signed chordal distance to the middle boundary point. The nearest treated and control locations remain roughly 0.31 and 0.42 km away.}
\end{frame}

\begin{frame}{Geographic estimates and uncertainty}
\begin{center}
\small\renewcommand{\arraystretch}{1.3}
\begin{tabular}{@{}lrrr@{}}\toprule
 & Estimate & Bias-corrected & Robust 95\% CI\\\midrule
Boundary point 1 & 2.87 & 3.63 & $[-6.40, 13.65]$\\
Boundary point 2 & -0.20 & -1.87 & $[-19.74, 16.01]$\\
Boundary point 3 & 4.78 & 4.71 & $[-19.64, 29.06]$\\
Whole-boundary distance & 4.95 & 6.37 & $[-8.34, 21.08]$\\
\bottomrule\end{tabular}
\end{center}

\begin{items}
\item Units: percentage points of turnout. Point-specific fits use signed chordal kilometres.
\item The robust intervals are wide. They allow substantively meaningful effects in both directions.
\item Distance to the whole boundary pools locations and changes the target.
\end{items}
\src{Module 6; Extensions, Section 5.2.3. The reported baseline does not model spatial dependence.}
\end{frame}

\begin{frame}{Extensions exercises and discussion}
\begin{items}
\item \textbf{B:} distinguish the Senate cutoff effect, window effect, and Fisherian null.
\item \textbf{C:} explain SPP eligibility effects, receipt LATEs, and window selection failure.
\item \textbf{D:} evaluate the probation tie rule, support gap, and outcome selection.
\item \textbf{E--F:} interpret pooling, boundary-specific effects, distance metrics, and geographic uncertainty.
\end{items}
\src{Worksheets B--F and Stata modules 2--6. Use selected modules in class and retain the others for a second lab.}
\end{frame}
